% \chapter*{Abstract}

% Effective field theories (EFTs) describe the low-energy effects of heavy new physics through the
% Wilson coefficients of higher-dimensional operators, as in the Standard Model Effective Field Theory
% (SMEFT) and the EFT of axion-like particles (ALPs), light pseudoscalars that bear on the strong CP
% problem, the flavor puzzle, and the nature of dark matter. Connecting experimental data with the underlying theory
% requires both the evolution of the Wilson coefficients with the energy, governed by their anomalous
% dimensions, and the region of parameter space compatible with the unitarity, causality, and locality
% of the scattering matrix. In this thesis, we address both questions via on-shell amplitude methods, which
% rely only on physical, gauge-invariant quantities and make manifest deep analytic and symmetry structures hidden in the
% Lagrangian formulation. We first derive on-shell formulae for the most general operator mixing up to two loops,
% including multiple operator insertions and leading mass effects through the Higgs low-energy
% theorem. We then compute the complete one-loop anomalous dimensions of the CP-violating ALP EFT,
% above and below the electroweak scale, and of the most general effective gauge theory up to dimension
% six, from which those of any specific model follow by group theory alone. Moreover, we show that
% helicity selection rules constrain the anomalous-dimension matrix at all loop orders, uncovering
% previously unknown zeros for operators of dimension five through eight, from two up to five loops. We
% then generalize partial-wave unitarity bounds to processes of arbitrary multiplicity and spin, and
% apply them to the quartic interactions of photons and gravitons, in combination with positivity
% bounds, and to $2\to3$ processes in the SMEFT. Furthermore, we derive unitarity and positivity
% bounds on the ALP EFT up to dimension eight, including weak-violating couplings to leptons, whose
% energy-enhanced effects in charged-meson and $W$-boson decays lead to bounds typically stronger than
% current experimental limits. Finally, we provide the first complete coupled-channel implementation
% of perturbative unitarity for the dimension-six SMEFT, supplemented by spinning sum rules and
% realistic flavor hypotheses, obtaining constraints competitive with experimental bounds above a few
% TeV.

\chapter*{Abstract}

Effective field theories (EFTs) describe the low-energy effects of heavy new physics through the Wilson coefficients of higher-dimensional operators.
Prominent examples are the Standard Model Effective Field Theory (SMEFT), the most general EFT built from the Standard Model fields alone, and the EFT of axion-like particles (ALPs), light pseudoscalars that bear on the strong CP problem, the flavor puzzle, and the nature of dark matter.
Connecting experimental data with the underlying theory requires both the evolution of the Wilson coefficients with the energy, governed by their anomalous dimensions, and the region of parameter space compatible with the unitarity, causality, and locality of the scattering matrix.
In this thesis, we address both aspects via on-shell amplitude methods, which rely only on physical, gauge-invariant quantities and make manifest deep analytic and symmetry structures hidden in the Lagrangian formulation. 
Regarding the anomalous dimensions, we first derive formulae that extract them from phase-space integrals of lower-loop amplitudes, valid for the most general operator mixing up to two loops, and include leading mass effects through the Higgs low-energy theorem while keeping the formalism massless.
We then compute the complete one-loop anomalous dimensions of the CP-violating ALP EFT, above and below the electroweak scale, and of the most general effective gauge theory up to dimension six, from which those of any specific model follow by group theory alone. Moreover, we show that helicity selection rules constrain the anomalous-dimension matrix at all loop orders, uncovering previously unknown zeros for operators of dimension five through eight. 
Concerning the theoretical constraints, we first generalize the formulation of partial-wave unitarity bounds and apply them to photon and graviton interactions.
We then derive unitarity bounds on the ALP EFT up to dimension eight, together with positivity bounds, which follow from analyticity and causality and fix the signs of combinations of Wilson coefficients. 
%For weak-violating couplings to leptons, whose effects in charged-meson and $W$-boson decays grow with the energy, the unitarity bounds are typically stronger than current experimental limits.
Finally, we provide the first complete coupled-channel implementation of perturbative unitarity for the dimension-six SMEFT and derive further bounds on four-fermion operators from spinning sum rules, dispersion relations that relate the signs of their coefficients to the spin of the heavy states. Under realistic hypotheses on the flavor structure, these constraints are competitive with experimental bounds above a few TeV.